\begin{abstract}
Research-agent systems now span automated machine-learning studies,
hypothesis generation, empirical software optimization, and laboratory-linked
analysis. A local success case therefore cannot support a broad claim of
autonomous-science performance. We study a narrower engineering question:
whether a persistent research state machine can bind a proposed successor to
a negative parent result, freeze an evaluation contract before outcome
access, retain failed rounds, and block claims that lack recorded evidence.
The study contains two equation-discovery rounds and a controlled workflow
matrix. Both equation-discovery candidates passed development screening and
failed their frozen unseen criteria. The workflow matrix contains
\TaskCount{} tasks, \IdempotencyRepeatCount{} deterministic repetitions,
three controller modes, and four ablations, for \CellCount{} recorded cells.
The repetitions test hash stability and do not add independent sampling
units. An additive reanalysis therefore treats task as the unit. The ten
within-task success differences are
\([0,0,0,1,0,0,1,1,1,1]\), with mean \TaskLevelGain{}, a
\BootstrapResamples{}-resample task bootstrap interval of
[\TaskLevelCILower{}, \TaskLevelCIUpper{}], and an exact two-sided sign-test
\(p\)-value of \SignTestTwoSided{}. The earlier 30-cell interval is retained
only as historical evidence for an internal gate. Exact cell replay and zero
unsupported claims remain useful engineering observations, but neither is an
external scientific effect. The result is a controlled demonstration of
evidence-state transitions and tamper detection. Independent tasks, common
budgets against external systems, interoperable provenance, and independent
human review remain required before any publication-facing performance
claim or submission decision.
\end{abstract}
